\subsection{Experimental Setup}
\label{sec:experimental_setup}

Seluruh eksperimen komputasi pada penelitian ini dijalankan pada infrastruktur perangkat keras yang mencakup prosesor Intel Core i9-13900H (14 inti, 20 utas), memori utama Random Access Memory (RAM) 32~GB DDR5 berkecepatan tinggi, serta NVIDIA GeForce RTX 4060 Laptop Graphics Processing Unit (GPU) dengan 8~GB Video Random Access Memory (VRAM). Alokasi beban komputasi menerapkan skema hibrida terkoordinasi demi efisiensi eksekusi pemrosesan data. Tahapan ekstraksi representasi laten Vision Transformer memanfaatkan akselerasi GPU NVIDIA melalui PyTorch, sedangkan pelatihan model dan penalaan hyperparameter validasi silang 5-fold untuk pengklasifikasi klasik dioptimalkan pada multi-inti Central Processing Unit (CPU) berbasis \texttt{scikit-learn}. Tumpukan perangkat lunak dibangun di atas Python 3.10 dengan pustaka PyTorch (\texttt{torch}) dan \texttt{scikit-learn} (\texttt{sklearn}) dalam lingkungan terisolasi untuk menjamin reproduktibilitas penuh seluruh prosedur eksperimen.
